\clearpage
\VerdierPhysicalPageAnchor{0211}%
\begingroup\small\setlength{\parskip}{0pt}\setlength{\emergencystretch}{1.5em}
\setlength{\abovedisplayskip}{4pt plus 2pt minus 2pt}%
\setlength{\belowdisplayskip}{4pt plus 2pt minus 2pt}%

\VerdierSemanticAnchor{ch3d-3-2-1-complex-x}
\noindent a) Le complexe $X$ dont le seul composant non nul est l'objet $X$
en degré zéro.

\VerdierSemanticAnchor{ch3d-3-2-1-complex-z}
\noindent b) Le complexe :
\[
Z^\bullet=\cdots\longrightarrow0\longrightarrow Y
\longrightarrow Z_{n-1}\longrightarrow\cdots\longrightarrow Z_0
\longrightarrow0\longrightarrow0\longrightarrow\cdots
\]
($Y$ est le composant de degré $-n$ et les $Z_i$ sont les composants de
degré $-i$).

\VerdierSemanticAnchor{ch3d-3-2-1-complex-y-shift}
\noindent c) Le complexe $Y[n]$ dont le seul composant non nul est l'objet
$Y$ en degré $-n$.

\VerdierSemanticAnchor{ch3d-3-2-1-morphisms}
Désignons alors par
$u(S^n(X,Y))\colon Z^\bullet\longrightarrow X$ le morphisme dont le seul
composant non nul est $u_0$ en degré $0$, et par
$v(S^n(X,Y))\colon Z^\bullet\longrightarrow Y[n]$ le morphisme dont le seul
composant non nul est $\operatorname{id}_Y$ en degré $-n$. Le morphisme
$u(S^n(X,Y))$ est un quasi-isomorphisme, d'où un morphisme dans
$D(\mathcal A)$ (avec les notations de (1.3.2)) :

\VerdierSemanticAnchor{ch3d-3-2-1-1}
\begin{equation}
\delta_{\mathcal A}^{n}(S^n(X,Y))=
[v(S^n(X,Y))][u(S^n(X,Y))]^{-1}\colon X\longrightarrow Y[n].
\tag{3.2.1.1}
\end{equation}

\VerdierSemanticAnchor{ch3d-3-2-1-2}
On définit ainsi une application :
\begin{equation}
\delta_{\mathcal A}^{n}\colon E_{\mathcal A}^{n}(X,Y)
\longrightarrow\operatorname{Ext}_{\mathcal A}^{n}(X,Y),
\tag{3.2.1.2}
\end{equation}
qu'on note simplement $\delta^n$ lorsqu'aucune confusion n'en résulte.

\VerdierSemanticAnchor{ch3d-3-2-2}
\medskip\noindent\textbf{Proposition 3.2.2.}\quad
\textit{a) L'application $\delta^n$ est surjective.}

\VerdierSemanticAnchor{ch3d-3-2-2-b-sequences}
\noindent\textit{b) Deux éléments :}
\[
S^n(X,Y)=0\longrightarrow Y\longrightarrow Z_{n-1}
\longrightarrow\cdots\longrightarrow Z_0\longrightarrow X
\longrightarrow0
\]
\noindent\textit{et :}
\[
{S'}^n(X,Y)=0\longrightarrow Y\longrightarrow Z'_{n-1}
\longrightarrow\cdots\longrightarrow Z'_0\longrightarrow X
\longrightarrow0
\]
\noindent\textit{ont la même image par $\delta^n$ si et seulement s'il
existe un élément :}
\[
{S''}^n(X,Y)=0\longrightarrow Y\longrightarrow Z''_{n-1}
\longrightarrow\cdots\longrightarrow Z''_0\longrightarrow X
\longrightarrow0
\]
\noindent\textit{et un diagramme commutatif :}

\VerdierSemanticAnchor{ch3d-3-2-2-b-diagram}
\begin{center}
\resizebox{\textwidth}{!}{%
$\begin{tikzcd}[ampersand replacement=\&,column sep=2.25em,row sep=2.5em]
0 \arrow[r] \& Y \arrow[r] \& Z_{n-1} \arrow[r] \& \cdots
  \arrow[r] \& Z_0 \arrow[r] \& X \arrow[r] \& 0 \\
0 \arrow[r] \& Y \arrow[r] \arrow[u,"{\operatorname{id}}"]
  \arrow[d,"{\operatorname{id}}"']
\& Z''_{n-1} \arrow[r] \arrow[u] \arrow[d]
\& \cdots \arrow[r]
\& Z''_0 \arrow[r] \arrow[u] \arrow[d]
\& X \arrow[r] \arrow[u,"{\operatorname{id}}"']
  \arrow[d,"{\operatorname{id}}"] \& 0 \\
0 \arrow[r] \& Y \arrow[r] \& Z'_{n-1} \arrow[r] \& \cdots
  \arrow[r] \& Z'_0 \arrow[r] \& X \arrow[r] \& 0
\end{tikzcd}$}
\end{center}
\endgroup
